% P05.2 one-vertex orthogonality classification, live integration.

Use the real scalar extension
$\mathrm{PA}_{\mathbf R}(\mathbf X)$ fixed in
\Cref{thm:frontier-polyhedral-one-dimensional-real-regularity}.  A
polyhedral convex combination is a finite barycentric combination
$x=\sum_j a_jx_j$ for which, on some common refinement, there are faces
$\tau$ and $\sigma_j$ with $x\in\tau$, $x_j\in\sigma_j$, and
$\tau\preceq\sigma_j$ for every $j$.  The envelope is
\[
 P_\gamma(u):=\sup\{\psi:\ \psi\leq\gamma+u,\quad
   \psi\in\mathrm{PC}(\mathbf X),\quad
   \psi\leq\gamma+O(1)\}.
\]
The preceding theorem proves, independently of the withdrawn supplied note, that
$P_\gamma(u)\in\mathrm{PA}_{\mathbf R}(\mathbf X)\cap
\mathrm{PC}(\mathbf X)$ on the same finite structure whenever
$u\in\mathrm{PA}_{\mathbf R,b}(\mathbf X)$.  For such a function, write
$\MA_{\mathrm{poly},\mathbf R}(f):=\Delta f$, where the
one-dimensional Laplacian uses the primitive integral-affine coordinate
and the balancing weight on each edge.  This is the real-linear extension
of the source's one-dimensional intersection formula
\cite[Definition~6.15 and Proposition~6.17]{BMPS26}.
We say that $(\mathbf X,\gamma)$ has real-polyhedral orthogonality if, for
every $u\in\mathrm{PA}_{\mathbf R,b}(\mathbf X)$,
\[
 \int_{\mathbf X}
 \bigl(P_\gamma(u)-(\gamma+u)\bigr)
 \MA_{\mathrm{poly},\mathbf R}(P_\gamma(u))=0.
\]

\begin{lemma}
\label{lem:frontier-polyhedral-one-vertex-convexity}
Let
\[
 \mathbf X=\bigcup_{i=1}^{r}\mathbf R_{\geq0}u_i
 \subset N_{\mathbf R}
\]
be the union of distinct rational rays with primitive generators $u_i$, put
\[
 K:=\left\{\lambda\in\mathbf R_{\geq0}^{r}:\
       \sum_{i=1}^{r}\lambda_i u_i=0\right\},
\]
and let $f$ be continuous and real-piecewise-affine on every ray.  If $s_i$
is the outgoing slope of $f$ at the origin along $u_i$, then $f$ is
polyhedrally convex if and only if every ray restriction is convex and
\[
                    \sum_i\lambda_i s_i\geq0
                    \quad\text{for every }\lambda\in K.
\]
If the ray $\mathbf R_{\geq0}u_i$ has balancing weight $w_i$, then
\[
             \MA_{\mathrm{poly},\mathbf R}(f)\{0\}=\sum_iw_i s_i.
\]
\end{lemma}

\begin{theorem}
\label{thm:frontier-polyhedral-one-vertex-orthogonality}
Let
$
 \mathbf X=\bigcup_{i=1}^{r}\mathbf R_{\geq0}u_i
$
be a one-dimensional balanced rational fan with distinct primitive ray
directions $u_i$ and positive rational weights $w_i$ such that
$
 \sum_iw_i u_i=0
$.
Let $\gamma\in\mathrm{PA}_{\mathbf R}(\mathbf X)$ be polyhedrally convex,
affine on every ray, with outgoing slopes $g_i>0$.  Then
$(\mathbf X,\gamma)$ has real-polyhedral orthogonality if and only if
\[
 \dim K=1,
 \qquad
 K=\left\{\lambda\in\mathbf R_{\geq0}^{r}:\
       \sum_i\lambda_i u_i=0\right\}.
\]
Equivalently,
$
 \dim_{\mathbf R}\operatorname{span}\{u_1,\ldots,u_r\}=r-1
$.
\end{theorem}

\begin{corollary}
\label{cor:frontier-polyhedral-one-vertex-bounded-shift}
In the setting of
\Cref{thm:frontier-polyhedral-one-vertex-orthogonality}, let $\gamma$ be
an arbitrary real-piecewise-affine polyhedrally convex reference positive at
infinity, without assuming that it is affine on each whole ray.  Thus
$\gamma$ is defined on a finite rational subdivision of the fan.  Then
$(\mathbf X,\gamma)$ has real-polyhedral orthogonality if and only if
$\dim K=1$.
\end{corollary}

\begin{example}
\label{ex:frontier-polyhedral-nonsaturated-tripod}
In $N=\mathbf Z^2$, take
\[
 u_1=(1,0),\qquad u_2=(1,2),\qquad u_3=(-3,-2),
 \qquad (w_1,w_2,w_3)=(2,1,1).
\]
The ray directions are primitive, their weighted sum is zero, and their
real span has rank two.  The lattice they generate has index two in $N$,
since the absolute values of the three pairwise determinants are $2,2,4$.
Thus the vertex is not polyhedrally smooth in the lattice sense, whereas
\Cref{cor:frontier-polyhedral-one-vertex-bounded-shift} gives
real-polyhedral orthogonality for every real-piecewise-affine
polyhedrally convex reference positive at infinity on this fan.
\end{example}
